% A2 — draft v0.1 (18/07/2026). Target: ECML-PKDD (alt.: IPM / Machine Learning j.)
% Estado: semente única (descritivo); campanha de 8 sementes NA FILA — converter
% achados 1-3 em inferenciais antes da submissão.
\documentclass[11pt]{article}
\usepackage[T1]{fontenc}
\usepackage[utf8]{inputenc}
\usepackage[english]{babel}
\usepackage{booktabs}
\usepackage{amsmath}
\usepackage[round]{natbib}
\usepackage[margin=2.8cm]{geometry}
\usepackage{microtype}
\usepackage{xcolor}
\newcommand{\todo}[1]{\textcolor{red}{[TODO: #1]}}

\title{The direction of self-evaluation bias in active learning depends on
the metric:\\ accuracy underestimates, macro-F1 overestimates}

\author{Gilsiley Henrique Darú$^{1}$ \and Gustavo Valentim Loch$^{1}$\\[2pt]
\normalsize $^{1}$Graduate Program in Numerical Methods in Engineering (PPGMNE),\\
\normalsize Federal University of Paraná (UFPR), Curitiba, Brazil}
\date{Draft v0.1 --- \today}

\begin{document}
\maketitle

\begin{abstract}
Practitioners of pool-based active learning routinely evaluate their model
on a held-out split of the labels they have collected --- because that is
all they have. It is folklore that this self-evaluation is biased, since
actively sampled data are not i.i.d.; but the direction and size of the
bias are rarely measured. We measure both, at scale, with a perfect-oracle
protocol that separates a labelable \emph{pool} (50{,}000 short product
descriptions, 649 classes, severe long tail) from a \emph{reserved
population} (181{,}490 instances) never touched by selection. At each step
the collected labels are split 80/20; the same model is scored on the
internal test split and on the population. Across five selection policies
and two classifiers we find the bias is large, systematic and
\emph{metric-dependent}: under uncertainty sampling, internal accuracy
\emph{underestimates} deployment accuracy by $17\pm1$ points across eight
seeds (the internal test inherits the hard cases), while internal macro-F1
\emph{overestimates} it by up to 34 points (the active sample
over-represents rare classes). A random-sampling control shows near-zero
gap, validating the instrument; entropy beats random in ceiling on 8 of 8
seeds (paired Wilcoxon, $p=0{,}0078$).
Two corollaries follow. First, \emph{more labels can hurt}: for a
logistic-loss classifier, training on the $\approx$15k actively selected
labels beats training on the full 50k pool by 15 macro-F1 points, because
the active sample acts as a balanced curriculum --- while a
per-class-normalized prototype classifier is immune. Second, an ablation
of a semantic-diversity selection policy shows its entire continuous-regime
gain comes from stratifying by predicted class, yielding a trivial,
strong, model-agnostic baseline. Self-evaluated metrics from active
learning runs should not be used for deployment decisions in either
direction; we quantify why, and what to do instead.
\end{abstract}

\noindent\textbf{Keywords:} active learning; sampling bias; evaluation;
class imbalance; macro-F1; text classification.

\section{Introduction}
\label{sec:intro}

Pool-based active learning (AL) buys labels selectively
\citep{Settles2012,Lewis1994}. Its success metric --- performance per label
--- is usually reported on a benchmark with abundant held-out gold data.
Deployed AL has no such luxury: the only labeled data available for
evaluation are the labels the loop itself collected, and the standard
practice is to hold out a fraction of them. The resulting sample is not
i.i.d.\ by construction: uncertainty sampling over-represents boundary
cases and rare classes \citep{Attenberg2010,Santos2016Vies}. That this
biases evaluation is known in principle; \emph{how much}, \emph{in which
direction}, and \emph{how the direction interacts with the metric} is
rarely quantified --- yet these are exactly the questions a practitioner
needs answered before trusting an internally computed number for a release
decision.

We answer them with a deliberately simple protocol (Section~\ref{sec:protocol}):
a perfect oracle (gold labels), a fixed labelable pool, and a large
reserved population that selection can never touch. The gap between
internal and population scores, tracked along the full labeling
trajectory, \emph{is} the self-evaluation bias --- and a random-sampling
control arm certifies that the gap is caused by selection, not by the
instrument.

Contributions: (i) the first (to our knowledge) large-scale, controlled
quantification of AL self-evaluation bias with 649 classes, showing it is
\textbf{bidirectional and metric-dependent} --- internal accuracy
underestimates deployment performance while internal macro-F1
overestimates it; (ii) a mechanistic account and empirical demonstration
of the \emph{less-is-more} effect --- full-pool labeling can be worse than
stopping at 30\% for loss functions dominated by the natural class
distribution --- including a classifier (per-class prototype) immune by
construction; (iii) an ablation showing a semantic-diversity continuous
policy owes its gain entirely to prediction-stratification, exposing a
trivial strong baseline; (iv) operational guidance: what internal metrics
can and cannot support, and a reserved-set release criterion.

\section{Related work}
\label{sec:related}

Sampling bias in AL is acknowledged since the field's foundations
\citep{Settles2012}; failure modes of uncertainty sampling under imbalance
are documented \citep{Attenberg2010}, and bias-aware AL has been studied
theoretically \citep{Santos2016Vies}. Density and cluster-based policies
\citep{Settles2008,Nguyen2004,Dasgupta2008} motivate the diversity
component we ablate. Metric behavior under imbalance (micro vs.\ macro
aggregation) is classical \citep{Sokolova2009}. What is missing in this
literature --- and what we contribute --- is a controlled, trajectory-level
measurement of the \emph{evaluation} bias (not the training bias) with its
metric dependence, at a class cardinality where the two metrics dissociate
dramatically. The renewed practical interest comes from LLM-oracle
pipelines \citep{Gilardi2023,Bayer2024}, where cheap labels make
full-trajectory experiments affordable and self-evaluation is the default
telemetry.

\section{Protocol: internal vs.\ population evaluation}
\label{sec:protocol}

Let $D$ be a deduplicated corpus, shuffled once with a fixed seed and split
into a \textbf{pool} $P$ ($|P|=50{,}000$), from which the learner may
request labels, and a \textbf{population} $Q$ ($|Q|=181{,}490$), reserved
in full. The oracle is perfect (gold labels), removing annotation noise as
a confounder. The loop starts from a random batch and iterates: train on
the collected set $L_t$, select the next batch of 500 from $P \setminus
L_t$, until the pool is exhausted. At every step, $L_t$ is split 80/20
(stratified when feasible, fixed seed); the model is trained on the 80\%
and scored twice with identical weights: on the \textbf{internal test}
(the 20\% --- what a practitioner without reserved data can measure) and on
the \textbf{population} $Q$ (deployment performance). The signed gap
$\Delta_m(t) = m_{\text{int}}(t) - m_{\text{ext}}(t)$ for metric $m$ is the
self-evaluation bias. Both scores use the same computation convention:
macro-F1 averages over the classes present in the respective reference
set, which is precisely the support asymmetry a practitioner faces ---
the internal test can only contain collected classes. Because the
internal test is small early in the trajectory (20\% of $|L_t|$),
single-step gaps are noisy; all reported biases are trajectory-level
patterns sustained over many consecutive steps, not single-point
readings.

\paragraph{Control arm.} Under random selection, $L_t$ is an i.i.d.\ sample
of $P$ (and $P$ of $D$), so $\Delta$ should vanish up to noise. A non-zero
$\Delta$ under random selection would indicate an instrument artifact
(split leakage, size effects); observing $\Delta \approx 0$ validates the
protocol, and every deviation under a non-random policy is then
attributable to selection.

\section{Experimental setup}
\label{sec:setup}

\textbf{Data and task.} 231{,}490 deduplicated short product descriptions
from Brazilian fiscal receipts (4--50 characters, uppercase,
abbreviations), 649 classes with at least 2 instances, severe long tail.
\textbf{Classifiers}, chosen as extremes of loss geometry: \emph{SGD} ---
logistic regression on binary TF--IDF via stochastic gradient descent,
loss dominated by the empirical class distribution; \emph{PVBin} --- a
per-class prototype model over binary term vectors
\citep{Daru2024Dissertacao}, normalized per class by construction.
\textbf{Policies}: uncertainty (entropy); random (control); a
semantic-diversity policy (density-representative selection over SBERT
embeddings with lexical-novelty tie-breaking, DRI-SL); the same policy
with groups given by the classifier's \emph{predicted classes} (DRI-SL-C);
and the ablation of the latter without lexical novelty --- equivalent to
\emph{stratified sampling by predicted class}. Batch 500, budget to pool
exhaustion. The entropy and random arms (the paper's central comparison)
were run over \textbf{eight seeds} that vary the initial draw, tie-breaks
and internal split while holding the pool fixed; results for these arms are
reported as mean$\pm$SD with a paired Wilcoxon test across seeds. The
diversity variants (DRI-SL, DRI-SL-C and its ablation) were run single-seed
and are reported descriptively.

\section{Results}
\label{sec:results}

\subsection{Selector comparison on the population}

\begin{table}[htb]
\centering
\caption{Population macro-F1 by policy and classifier. Entropy and random
over 8 seeds (mean$\pm$SD); diversity variants single-seed (descriptive).
Saturation: smallest $|L|$ reaching 95\% of the arm's own ceiling.}
\label{tab:selectors}
\small
\begin{tabular}{llrrrr}
\toprule
\textbf{Clf.} & \textbf{Policy} & \textbf{Ceiling} & \textbf{Saturation} &
\textbf{F1@10k} & \textbf{F1@20k} \\
\midrule
SGD & Entropy      & \textbf{0.592$\pm$0.004} & \textbf{9{,}060$\pm$560} & 0.565 & 0.574 \\
SGD & Stratif.\ by prediction & 0.555 & 10{,}000 & 0.533 & 0.543 \\
SGD & DRI-SL-C     & 0.491 & 15{,}500 & 0.412 & 0.486 \\
SGD & Random       & 0.455$\pm$0.004 & 15{,}500$\pm$1{,}100 & 0.391 & 0.449 \\
SGD & DRI-SL       & 0.441 & 41{,}500 & 0.310 & 0.365 \\
\midrule
PVBin & Entropy    & 0.531$\pm$0.002 & 19{,}310$\pm$700 & 0.438 & 0.509 \\
PVBin & Stratif.\ by prediction & 0.528 & 18{,}000 & 0.444 & 0.498 \\
PVBin & DRI-SL-C   & 0.525 & 39{,}500 & 0.349 & 0.453 \\
PVBin & Random     & 0.528$\pm$0.004 & 40{,}630$\pm$1{,}410 & 0.348 & 0.423 \\
PVBin & DRI-SL     & 0.527 & 45{,}500 & 0.284 & 0.356 \\
\bottomrule
\end{tabular}
\end{table}

Entropy dominates at scale (Table~\ref{tab:selectors}), and the eight-seed
campaign shows the dominance is not a lucky draw: for SGD it saturates at
$9{,}060\pm560$ labels (range 8{,}500--10{,}000) against $15{,}500\pm1{,}100$
for random, and the entropy ceiling exceeds random's in \textbf{8 of 8
seeds} (paired Wilcoxon, the maximum significance reachable with eight
pairs, $p=0{,}0078$); for PVBin the entropy edge in ceiling is small and
not significant (6/8 seeds), but its saturation ($19{,}310\pm700$) is
still less than half of random's ($40{,}630\pm1{,}410$). The
semantic-diversity policy (DRI-SL), designed for the zero-label cold
start, is \emph{worse than random} as a continuous policy; guiding its
groups by predicted classes (DRI-SL-C) repairs it (ceiling
$0.441 \to 0.491$ for SGD, saturation $2.7\times$ earlier).

\subsection{The ablation: stratification is the active ingredient}

Removing the lexical-novelty term from DRI-SL-C --- leaving plain
stratified sampling by predicted class --- \emph{improves} it on both
classifiers (SGD: ceiling 0.555 vs.\ 0.491, saturation 10k vs.\ 15.5k;
PVBin: saturation 18k vs.\ 39.5k, tying entropy). The diversity heuristic,
valuable at cold start where no model exists, is counterproductive in the
continuous regime: it privileges lexically atypical exemplars when the
macro metric pays for balanced class coverage. The practical corollary is
a baseline: \emph{predict, stratify, sample} --- model-agnostic, no
embeddings, no tuning --- that captures most of the uncertainty-sampling
gain and nearly all of the diversity-policy gain.

\subsection{The bias is bidirectional and metric-dependent}

The control behaves: under random selection the internal--external
accuracy gap is $\Delta \approx 0$ along the whole trajectory. Under
entropy selection, internal \textbf{accuracy underestimates} the
population value substantially and \emph{robustly across seeds} --- at
$|L|=10{,}000$ the gap is $-17{,}1\pm1{,}0$ points over the eight seeds
(tight SD): the internal test inherits the boundary-heavy, hard-case
distribution that uncertainty sampling selects. Internal \textbf{macro-F1
overestimates} the population value for \emph{every} non-random policy ---
by up to \textbf{$+34$ points} early in the DRI-SL trajectory (single
seed): the active sample gives rare classes internal test support far
above their population share, so per-class scores are computed on
flattering supports. Neither direction is usable as a deployment
estimate, and the two errors do not cancel: a practitioner reading
internal accuracy will under-deploy (the model is better than it looks);
one reading internal macro-F1 will over-promise tail quality (it is worse
than it looks).

\subsection{More labels can hurt: the active sample as curriculum}

For SGD under entropy selection, population macro-F1 peaks at
$0.59$ around 15k labels and \emph{falls} to $0.44$ when the entire pool
is labeled: the final 35k labels, drawn increasingly from the natural
(imbalanced) distribution, drag the logistic loss toward the head classes.
The actively selected subset acts as a balanced curriculum --- and
labeling everything undoes it. PVBin, whose per-class prototypes normalize
away the class distribution, is immune (ceiling $\approx$ flat across
policies). Beyond waste, additional labels can be \emph{harmful} to macro
metrics for distribution-sensitive losses; remedies are early stopping by
stagnation, reweighting, or per-class-normalized models.

\section{Operational guidance}
\label{sec:guidance}

(i) Maintain a reserved evaluation set from day one --- the bias cannot be
corrected post hoc from internal data alone, since its size and sign
depend on the (policy, metric, model) triple. (ii) Use internal metrics
only for \emph{relative} decisions inside one run (stagnation-based
stopping), never for absolute release decisions. (iii) If a reserved set
is impossible, prefer internal accuracy over internal macro-F1 for
go/no-go: it errs conservative. (iv) For macro objectives with
distribution-sensitive losses, treat the labeled-set composition as part
of the model: stopping earlier, subsampling or reweighting may beat
labeling more.

\section{Limitations}
\label{sec:limits}

Single dataset and domain (short retail text; the mechanism --- support
inflation of rare classes --- is general, but effect sizes are not); the
central entropy-vs-random comparison is inferential over eight seeds, but
the diversity-variant arms (DRI-SL, DRI-SL-C, ablation) remain single-seed
and descriptive, so their between-arm orderings should be read as
indicative pending a seed campaign of their own; perfect oracle by design
(noisy oracles would add a second bias source; our companion work on LLM
oracles addresses it separately, keeping this paper's question clean).

\section{Conclusion}
\label{sec:conclusion}

Self-evaluation on actively collected data is not merely noisy --- it is
systematically wrong in a direction that depends on the metric: accuracy
reads low, macro-F1 reads high, by double-digit margins at 649-class
granularity. A random control certifies the effect is caused by selection.
The same experiments yield a practical bonus (stratified-by-prediction as
a trivial strong policy) and a warning (labeling the full pool can damage
macro performance for distribution-sensitive losses). All trajectories,
code and per-step artifacts are released.

\paragraph{Artifacts.} \texttt{github.com/GHDaru/activelearning}
(\texttt{experiments/e6population}).

\bibliographystyle{plainnat}
\bibliography{refs}

\end{document}
